\chapter{Empirical assessment of the control framework}
\section{Experimental protocol and evaluation setup}
% \subsection{Physical and numerical model}
% \label{subsec:experimental_protocol_setup}

% All numerical experiments in this thesis are carried out within the same effective control framework introduced in the previous chapter, but instantiated through a concrete and uniform simulation pipeline. The controlled system is modeled as a two-qubit gmon architecture truncated to three local levels per mode, so that the full physical Hilbert space has dimension \(3^2=9\), while the computational subspace remains four-dimensional. The underlying dynamics are generated by the effective rotating-wave Hamiltonian already introduced in Eq.~\eqref{eq:gmon_rwa_bosonic_hamiltonian}, and the target operations belong to the family \(N(\alpha,\alpha,\gamma)\) defined in Eq.~\eqref{eq:target_gate_family}. In this way, the numerical section does not introduce a new physical model, but rather provides a consistent implementation of the same multilevel control problem discussed theoretically above. Unless otherwise stated, the simulations use the default truncation and control setting adopted in the codebase, namely local anharmonicity centered around a fixed baseline value \(\eta_{\mathrm{base}}=-200~\mathrm{MHz}\), a control amplitude scale bounded by \(20~\mathrm{MHz}\), and a finite-bandwidth filtering stage with characteristic bandwidth \(10~\mathrm{MHz}\), so as to remain compatible with the physical assumptions underlying the UFO framework.

% The methods compared in this section are chosen so as to isolate the different contributions of the framework. First, a direct gradient-based baseline is implemented through an Adam optimizer acting directly on the control trajectory; this provides a non-RL reference that optimizes the same physical problem without policy learning. Second, a TRPO-based controller is trained in a nominal environment, that is, without stochastic perturbations during training. Third, the same TRPO controller is trained in the stochastic setting described in the previous chapter, producing the noise-optimized RL solution that is most closely aligned with the methodology of the original paper. Finally, for runtime comparisons on the gate family \(N(\alpha,\alpha,\pi/2)\), the analog-control solutions are compared against the optimal gate-synthesis baseline discussed in Section~\ref{subsec:baselines_comparison_strategy}. This choice of baselines is deliberate: the Adam baseline probes whether policy learning is advantageous over direct pulse optimization, the nominal TRPO baseline isolates the effect of noisy training, and the gate-synthesis baseline probes whether direct analog control provides a hardware-efficiency advantage over compiled circuit implementations.

% The evaluation metrics are organized according to the distinct physical questions addressed by the experiments. For a deterministic control plan \(\mathcal{P}\), the nominal assessment uses the final projected unitary \(K_{\mathcal{P}}\), the corresponding gate fidelity with respect to the target operation, the TSWT-based leakage estimate, the total runtime, and the complete UFO cost. In compact form, the nominal fidelity can be derived from \ref{eq:gate_fidelity_unitary}, and given the computational subspace dimension \(d=4\), is computed as
% \begin{equation}
% F_{\mathrm{nom}}(\mathcal{P})
% =
% \frac{1}{16}
% \left|
% \mathrm{Tr}\!\left(
% K_{\mathcal{P}}^{\dagger}U_{\mathrm{target}}
% \right)
% \right|^{2},
% \label{eq:nominal_fidelity_validation}
% \end{equation}
% while the nominal control cost is evaluated through the same UFO functional already introduced in Eq.~\eqref{eq:ufo_cost_function}. 
% For robustness studies, deterministic evaluation is supplemented by Monte Carlo sampling of noisy control realizations. 
% In that setting, three quantities are reported: the average channel fidelity estimator \ref{eq:average_gate_fidelity_def} used in the implementation, the average gate fidelity over noisy realizations \ref{eq:TODO}, and the fidelity variance \ref{eq:TODO}. Denoting a noisy realization of the control trajectory by \(\mathcal{P}_{\xi}\), one may write schematically
% \begin{equation}
% \overline{F}_{\mathrm{gate}}
% =
% \mathbb{E}_{\xi}
% \!\left[
% F_{\mathrm{gate}}(\mathcal{P}_{\xi})
% \right],
% \qquad
% \mathrm{Var}_{\xi}(F_{\mathrm{gate}})
% =
% \mathbb{E}_{\xi}
% \!\left[
% \bigl(F_{\mathrm{gate}}(\mathcal{P}_{\xi})-\overline{F}_{\mathrm{gate}}\bigr)^2
% \right],
% \label{eq:robustness_metrics_validation}
% \end{equation}
% so that the numerical study captures not only average performance under perturbations, but also the shot-to-shot stability of each control strategy. Runtime is then treated as a separate metric when comparing analog solutions against the gate-synthesis baseline.

% All methods are evaluated under a common protocol designed to make the comparisons as fair and interpretable as possible. The RL environment discretizes the pulse into steps of duration \(\Delta t\), propagates the full \(9\times 9\) unitary evolution, and computes rewards from the current UFO cost, while deterministic post-training evaluation is always performed in the nominal setting unless robustness is explicitly being tested. In the robustness analysis, noisy realizations are generated by adding Gaussian perturbations to the same physical control channels that appear in the Hamiltonian, together with fluctuations in the anharmonicity parameter, so that the evaluation noise model remains aligned with the training environment and with the assumptions of the paper. 
% The nominal and noisy metrics are then computed from saved control plans rather than from transient training trajectories, ensuring that the reported results refer to fully specified pulse solutions. 
% % This separation between training and evaluation is important for the logic of the thesis: the former determines how the controller is learned, whereas the latter determines how its physical quality is judged. 
% In the sections that follow, this common protocol will be used first to validate the framework on representative targets and then to compare its robustness and runtime behavior across the experiments of interest.


% \paragraph{Light hyperparameter sensitivity study.}
% \label{par:light_hyperparameter_sensitivity}

% Before running the main validation experiments, a limited hyperparameter sensitivity study was performed in order to select a stable TRPO configuration for the subsequent analysis. The purpose of this search was not to perform a large-scale optimization over the full algorithmic design space, but rather to identify a reasonable and reproducible operating point for a representative target gate. In particular, the search was carried out on the target \(N(\alpha,\alpha,\gamma)\) with
% \begin{equation}
% \alpha = 2.2,
% \qquad
% \gamma = \pi/2,
% \label{eq:representative_target_tuning}
% \end{equation}
% and explored a small grid over the most influential TRPO hyperparameters, namely the trust-region radius \(\texttt{max\_kl}\), the stopping threshold \(\texttt{termination\_cost}\), the initial policy exploration scale \(\texttt{init\_log\_std}\), and the temporal discretization / horizon pair \((dt, T_{\max})\). The search was intentionally kept lightweight: it covered \(9\) TRPO configurations and \(27\) total runs, corresponding to \(3\) random seeds per configuration. This limited design is consistent with the methodological role of the study, namely parameter selection rather than a standalone scientific result. 
% % :contentReference[oaicite:0]{index=0}

% The configuration retained for the main experiments was the one with the lowest aggregated primary metric across seeds. The selected setting corresponds to
% \begin{equation}
% \texttt{max\_kl}=0.005,
% \qquad
% \texttt{termination\_cost}=0.15,
% \qquad
% \texttt{init\_log\_std}=-1.0,
% \qquad
% dt=2.0~\mathrm{ns},
% \qquad
% T_{\max}=180.0~\mathrm{ns},
% \label{eq:selected_trpo_hyperparameters}
% \end{equation}
% for the representative target in Eq.~\eqref{eq:representative_target_tuning}. Averaged over the three seeds used in the search, this configuration achieved a best evaluation cost of \(0.3714\), a best evaluation fidelity of \(0.9941\), and a best evaluation leakage of \(3.91\times 10^{-4}\), with negligible dispersion in runtime since all runs converged to \(180~\mathrm{ns}\). These values were considered sufficiently stable to justify adopting the configuration in Eq.~\eqref{eq:selected_trpo_hyperparameters} for the main validation experiments reported below. 

% % :contentReference[oaicite:1]{index=1}

% \subsection{Single-target control optimization}
% \label{subsec:single_target_control_optimization}

% As a first validation step, the framework was tested on a single representative target gate from the family \(N(\alpha,\alpha,\gamma)\). The choice adopted throughout this subsection is
% \begin{equation}
% \alpha = 2.2,
% \qquad
% \gamma = \pi/2,
% \label{eq:representative_single_target}
% \end{equation}
% which was also used as the reference target in the light hyperparameter sensitivity study described in Section~\ref{subsec:experimental_protocol_setup}. This choice is convenient for several reasons. First, it provides a nontrivial two-qubit target that is sufficiently far from the identity to make the control problem genuinely meaningful. Second, it belongs to the subfamily with \(\gamma=\pi/2\), which is also the one employed later in the runtime analysis, thereby ensuring continuity between the single-target study and the broader family-level investigation. Finally, the selected target was found to admit stable and reproducible optimization behavior across multiple seeds during the preliminary parameter search, making it a suitable benchmark for comparing the Adam baseline, nominal TRPO, and noise-optimized TRPO under a common and well-controlled setting. 

% % :contentReference[oaicite:1]{index=1}

% The first reference solution for the target in Eq.~\eqref{eq:representative_single_target} was obtained through direct gradient-based optimization of the control trajectory using the Adam baseline. In this setting, the control problem is optimized independently for a discrete set of candidate horizons, and the best solution is then selected according to the nominal UFO cost. For the present target, the horizon sweep
% \begin{equation}
% T \in \{60,90,120,150,180,210,240\}\ \mathrm{ns}
% \label{eq:adam_horizon_grid}
% \end{equation}
% identified
% \begin{equation}
% T^{\star}_{\mathrm{Adam}} = 90~\mathrm{ns}
% \label{eq:adam_best_horizon}
% \end{equation}
% as the best-performing duration, with objective cost
% \begin{equation}
% C_{\mathrm{Adam}}^{\star} = 0.15468,
% \label{eq:adam_best_cost}
% \end{equation}
% nominal fidelity
% \begin{equation}
% F_{\mathrm{Adam}}^{\mathrm{nom}} = 0.999651,
% \label{eq:adam_nominal_fidelity}
% \end{equation}
% and nominal leakage
% \begin{equation}
% L_{\mathrm{Adam}}^{\mathrm{nom}} = 9.59\times 10^{-5}.
% \label{eq:adam_nominal_leakage}
% \end{equation}
% The corresponding runtime is, by construction,
% \begin{equation}
% T_{\mathrm{Adam}}^{\mathrm{nom}} = 90~\mathrm{ns}.
% \label{eq:adam_nominal_runtime}
% \end{equation}
% The horizon search is also informative from a qualitative viewpoint: the best cost is attained at an intermediate duration, while both shorter and longer horizons lead to worse objective values, with the cost increasing from \(0.3413\) at \(60~\mathrm{ns}\) to \(0.4005\) at \(240~\mathrm{ns}\). This behavior is physically consistent with the multi-objective structure of the UFO cost. Extremely short horizons make the control problem more aggressive and therefore harder to solve accurately, whereas overly long horizons are penalized by the runtime term and offer no compensating gain in fidelity. The Adam baseline therefore provides a strong nominal reference solution for this target: it already achieves very high gate fidelity and low leakage, while also revealing that the optimal operating regime is not at the shortest admissible time, but at a moderate horizon where the tradeoff between accuracy, leakage suppression, and runtime is best balanced.

% \subsection{Experimental protocol and evaluation setup}
\label{subsec:experimental_protocol_setup}

All numerical experiments in this thesis are performed within the same effective control framework introduced in the previous section. The physical system is modeled as a pair of weakly anharmonic superconducting modes truncated to three local levels each, so that the full Hilbert space has dimension
\(d_{\mathrm{phys}} = 3^2 = 9\),
% \begin{equation}
% d_{\mathrm{phys}} = 3^2 = 9,
% \label{eq:validation_physical_dimension}
% \end{equation}
while the computational two-qubit subspace retains dimension
\(d_{\mathrm{comp}} = 4\).
% \begin{equation}
% d_{\mathrm{comp}} = 4.
% \label{eq:validation_computational_dimension}
% \end{equation}
The dynamics are generated by the effective gmon Hamiltonian introduced in Section~\ref{subsec:gmon_effective_hamiltonian}, and are simulated in discrete time with timestep \(\Delta t\) over a finite control horizon \(T\). At each time step, the control action specifies the seven-dimensional control vector
\begin{equation}
\bm{u}(t)
=
\bigl(
g(t),\,
\delta_1(t),\delta_2(t),\,
f_1(t),f_2(t),\,
\phi_1(t),\phi_2(t)
\bigr),
\label{eq:validation_control_vector}
\end{equation}
consisting of a tunable coupling pulse, two local detuning controls, two microwave amplitudes, and two microwave phases. These controls are passed through a finite-bandwidth filter before being applied to the Hamiltonian, so that the simulated pulses remain consistent with the experimentally motivated smooth-control assumptions built into the framework. For every experiment, the target operation is taken from the gate family
\begin{equation}
N(\alpha,\alpha,\gamma)
=
\exp\!\left[
i\left(
\alpha\,\sigma_1^x\sigma_2^x
+
\alpha\,\sigma_1^y\sigma_2^y
+
\gamma\,\sigma_1^z\sigma_2^z
\right)
\right],
\label{eq:validation_target_family}
\end{equation}
and control quality is assessed through the same UFO cost structure introduced previously, namely a weighted combination of gate infidelity, leakage penalty, boundary penalty, and runtime. This uniform physical and numerical setup is essential for the interpretation of the experiments, because it ensures that all methods are compared within the same multilevel model, under the same control constraints, and against the same class of target gates.

TODO: Aggiungere un diagramma parlante

\subsection{Compared methods}
Within this common framework, four optimization and comparison schemes are considered. 
\begin{itemize}
    \item The first is a direct gradient-based baseline, implemented through the Adam optimizer, which searches directly over the filtered control trajectory for a fixed gate horizon. 
    \item The second is a policy-gradient reinforcement-learning controller trained in the nominal environment, that is, without stochastic perturbations applied during training. 
    \item The third is the noise-optimized TRPO controller, which is trained in a stochastic environment where Gaussian perturbations are injected into the control amplitudes and into the anharmonicity-related parameter entering the Hamiltonian. 
    \item The fourth comparison is not a pulse-level optimizer but a circuit-level reference, namely the optimal gate-synthesis baseline used to benchmark runtime on the family \(N(\alpha,\alpha,\gamma)\). 
\end{itemize}
These four methods play different methodological roles. The Adam baseline tests whether a standard first-order optimization of the raw control variables is sufficient to obtain high-quality pulses. The nominal TRPO controller tests whether reinforcement learning is already competitive in the idealized deterministic setting. The noise-optimized TRPO controller tests whether robustness can be learned directly from stochastic training. Finally, the gate-synthesis baseline tests whether direct analog control offers a runtime advantage over a discrete compiled implementation of the same target gate. Taken together, these methods define a comparison protocol that is not centered on a single optimizer, but on a set of scientifically distinct questions concerning accuracy, robustness, and hardware efficiency.


The performance of each control method is evaluated through a set of complementary metrics designed to probe different aspects of the control problem. At the most basic level, the implemented gate is assessed through its fidelity with respect to the target unitary. For a final physical propagator \(U(T)\), projected onto the computational subspace, the gate fidelity is computed as
\begin{equation}
F_{\mathrm{gate}}
=
\frac{1}{16}
\left|
\mathrm{Tr}\!\left(
U_{\mathrm{proj}}^{\dagger}(T)\,U_{\mathrm{target}}
\right)
\right|^2,
\label{eq:validation_gate_fidelity}
\end{equation}
where \(U_{\mathrm{target}}=N(\alpha,\alpha,\gamma)\). This quantity is reported together with the total UFO cost, which combines fidelity, leakage, boundary penalties, and runtime into a single scalar objective. Since the framework is explicitly multilevel, nominal performance is not characterized by fidelity alone. For every optimized control plan, the simulation also returns the TSWT-based leakage estimate \(L_{\mathrm{tot}}\), its decomposition into boundary and integral contributions, the boundary penalty term, and the final runtime in nanoseconds. These quantities allow one to distinguish between a control that is merely accurate on the logical subspace and one that is also physically well behaved in the full truncated Hilbert space. In addition to nominal performance, robustness is evaluated statistically under sampled noisy control realizations. Given a control plan \(\tilde{\bm{c}}\) and a Gaussian noise level \(\sigma_{\mathrm{noise}}\), the key robustness diagnostics are the average fidelity
\begin{equation}
\overline{F}
=
\mathbb{E}_{\delta\tilde{\bm{c}}}
\Bigl[
F\bigl(U(\tilde{\bm{c}}+\delta\tilde{\bm{c}})\bigr)
\Bigr],
\label{eq:validation_average_fidelity}
\end{equation}
the corresponding average gate fidelity, and the fidelity variance
\begin{equation}
\mathrm{Var}(F)
=
\mathbb{E}_{\delta\tilde{\bm{c}}}
\left[
\bigl(F-\overline{F}\bigr)^2
\right].
\label{eq:validation_fidelity_variance}
\end{equation}
These metrics are essential because they distinguish mean robustness from shot-to-shot stability: a control scheme with high average fidelity but large variance remains experimentally fragile. Finally, for the runtime study, the relevant metric is the total control duration required to realize a target gate, compared against the optimal discrete gate-synthesis baseline. The resulting evaluation protocol is therefore explicitly multidimensional: it measures logical accuracy, leakage suppression, robustness under stochastic perturbations, and hardware efficiency rather than relying on a single scalar score.

All experiments follow a common procedural logic. For a given target gate \(N(\alpha,\alpha,\gamma)\), a control method is first trained or optimized under its own learning rule until either a prescribed iteration budget is exhausted or a termination criterion on the UFO cost is satisfied. The resulting best control plan is then stored and re-evaluated in a deterministic nominal setting in order to extract its final cost, fidelity, leakage, and runtime under identical physical assumptions. For robustness studies, the same saved control plan is subsequently subjected to repeated noisy rollouts, in which Gaussian perturbations are independently sampled and injected into the control channels and into the anharmonicity-related parameter of the Hamiltonian. This Monte Carlo evaluation yields the empirical estimates entering Eqs.~\eqref{eq:validation_average_fidelity} and \eqref{eq:validation_fidelity_variance}. In the runtime study, by contrast, the focus is not on statistical robustness but on how the optimized gate duration varies across the target family, especially along the subfamily \(N(\alpha,\alpha,\pi/2)\), where analog control can be meaningfully compared with optimal gate synthesis. The same timestep, control bounds, bandwidth assumptions, and physical truncation are retained throughout, so that differences between methods can be attributed to the optimization strategy rather than to inconsistencies in the simulation model. This common evaluation protocol is particularly important for the interpretation of the results: it ensures that the subsequent comparisons are methodologically fair, while also making the distinction between nominal accuracy, noisy robustness, and runtime efficiency completely explicit.
